\documentclass[10pt,a4paper]{amsart}
\usepackage{amsmath,amssymb,amsthm,graphicx}
% Retained minimal preamble for the Stacks Project chapter more-algebra.tex (More on Algebra).
% The chapter file pinned-original-source.tex carries only \input{preamble} on its first line, so no
% preamble text lives inside the chapter itself.  The Stacks Project uses one single shared preamble for
% every chapter; the snapshot stages it identically as sources/stacks-chow/preamble.tex and
% sources/stacks-etale-cohomology/preamble.tex.  The block below is copied mechanically, line for line,
% from sources/stacks-chow/preamble.tex lines 14-18 (the Xy-pic setup) and 154-217 (the native theorem
% declarations together with the mathematical macros).  Nothing was added, removed or redefined; the
% wrapper class and the amsmath/amssymb/amsthm/graphicx base packages are supplied by the assembler.
% Chapter: more-algebra.tex -- More on Algebra.
% For commutative diagrams we use Xy-pic
\usepackage[all]{xy}

% We use 2cell for 2-commutative diagrams.
\xyoption{2cell}

% Theorem environments.
%
\theoremstyle{plain}
\newtheorem{theorem}[subsection]{Theorem}
\newtheorem{proposition}[subsection]{Proposition}
\newtheorem{lemma}[subsection]{Lemma}

\theoremstyle{definition}
\newtheorem{definition}[subsection]{Definition}
\newtheorem{example}[subsection]{Example}
\newtheorem{exercise}[subsection]{Exercise}
\newtheorem{situation}[subsection]{Situation}

\theoremstyle{remark}
\newtheorem{remark}[subsection]{Remark}
\newtheorem{remarks}[subsection]{Remarks}

\numberwithin{equation}{subsection}

% Macros
%
\def\lim{\mathop{\mathrm{lim}}\nolimits}
\def\colim{\mathop{\mathrm{colim}}\nolimits}
\def\Spec{\mathop{\mathrm{Spec}}}
\def\Hom{\mathop{\mathrm{Hom}}\nolimits}
\def\Ext{\mathop{\mathrm{Ext}}\nolimits}
\def\SheafHom{\mathop{\mathcal{H}\!\mathit{om}}\nolimits}
\def\SheafExt{\mathop{\mathcal{E}\!\mathit{xt}}\nolimits}
\def\Sch{\mathit{Sch}}
\def\Mor{\mathop{\mathrm{Mor}}\nolimits}
\def\Ob{\mathop{\mathrm{Ob}}\nolimits}
\def\Sh{\mathop{\mathit{Sh}}\nolimits}
\def\NL{\mathop{N\!L}\nolimits}
\def\CH{\mathop{\mathrm{CH}}\nolimits}
\def\proetale{{pro\text{-}\acute{e}tale}}
\def\etale{{\acute{e}tale}}
\def\QCoh{\mathit{QCoh}}
\def\Ker{\mathop{\mathrm{Ker}}}
\def\Im{\mathop{\mathrm{Im}}}
\def\Coker{\mathop{\mathrm{Coker}}}
\def\Coim{\mathop{\mathrm{Coim}}}

% Boxtimes
%
\DeclareMathSymbol{\boxtimes}{\mathbin}{AMSa}{"02}

%
% Macros for moduli stacks/spaces
%
\def\QCohstack{\mathcal{QC}\!\mathit{oh}}
\def\Cohstack{\mathcal{C}\!\mathit{oh}}
\def\Spacesstack{\mathcal{S}\!\mathit{paces}}
\def\Quotfunctor{\mathrm{Quot}}
\def\Hilbfunctor{\mathrm{Hilb}}
\def\Curvesstack{\mathcal{C}\!\mathit{urves}}
\def\Polarizedstack{\mathcal{P}\!\mathit{olarized}}
\def\Complexesstack{\mathcal{C}\!\mathit{omplexes}}
% \Pic is the operator that assigns to X its picard group, usage \Pic(X)
% \Picardstack_{X/B} denotes the Picard stack of X over B
% \Picardfunctor_{X/B} denotes the Picard functor of X over B
\def\Pic{\mathop{\mathrm{Pic}}\nolimits}
\def\Picardstack{\mathcal{P}\!\mathit{ic}}
\def\Picardfunctor{\mathrm{Pic}}
\def\Deformationcategory{\mathcal{D}\!\mathit{ef}}
\usepackage{cleveref}
\crefname{originallemma}{original lemma}{original lemmas}
\Crefname{originallemma}{Original lemma}{Original lemmas}
\makeatletter
\newlabel{more-algebra-lemma-modules}{{0D2F}{source}{}{sourcepointer-1}{}}
\newlabel{more-algebra-lemma-modules@cref}{{[originallemma][0][]0D2F}{[1][0][]source}{}{}{}}
\makeatother
\crefname{originallemma}{original lemma}{original lemmas}
\Crefname{originallemma}{Original lemma}{Original lemmas}
\makeatletter
\newlabel{equation-modules}{{0D2E}{source}{}{sourcepointer-2}{}}
\newlabel{equation-modules@cref}{{[originallemma][0][]0D2E}{[1][0][]source}{}{}{}}
\makeatother
\crefname{originallemma}{original lemma}{original lemmas}
\Crefname{originallemma}{Original lemma}{Original lemmas}
\makeatletter
\newlabel{algebra-lemma-adjoint-tensor-restrict}{{05DQ}{source}{}{sourcepointer-3}{}}
\newlabel{algebra-lemma-adjoint-tensor-restrict@cref}{{[originallemma][0][]05DQ}{[1][0][]source}{}{}{}}
\makeatother
\title{The right adjoint of the base change functor over a fibre product of rings}
\author{The Stacks Project Authors}
\date{stacks-project:more-algebra.tex:local-snapshot-20261012; GNU FDL 1.2 or later}
\begin{document}
\maketitle
\setcounter{subsection}{1}
\noindent\textit{Editorial note.} All delivered mathematical text is the author's own, copied from the pinned Stacks Project chapter file more-algebra.tex; no statement or proof step was rewritten, reordered or removed.\par\medskip
\noindent\textit{Editorial note.} The chapter file pinned-original-source.tex carries only \textbackslash{}input\{preamble\} on its first line, so no preamble text lives inside the chapter. The retained file sources/stacks-more-algebra/retained-preamble.tex is a mechanical, line-for-line copy of the project's single shared preamble as staged in the snapshot (Xy-pic setup plus theorem declarations and mathematical macros). Nothing was added, removed or redefined, and the excerpt serves only as preamble, never as a span coordinate.\par\medskip
\noindent\textit{Editorial note.} Both counted spans start at the \textbackslash{}begin line of their environment, so the statement and the author's proof are delivered whole; surrounding prose of the section is deliberately outside the counted spans and is carried as standard prerequisites instead.\par\medskip
\noindent\textit{Editorial note.} The citation fields (source\_id, source\_version, authors, source\_url, license.name, source\_title) are identical to those of every other card this lane has produced from this chapter, so the batch selector treats them as one source.\par\medskip
\noindent\textit{Allowed standard prerequisites (editorial).}\begin{itemize}
\item The functor from modules over a fibre product of rings to pairs of modules appearing in the statement. the author refers to the displayed equation equation-modules, and the surrounding section, source lines 400-461. Supplies the functor whose right adjoint is constructed.
\item The tensor-restriction adjunction identifying maps out of a base changed module with maps out of the original one. the author cites Algebra, Lemma algebra-lemma-adjoint-tensor-restrict at proof lines 495-497. Rewrites the right hand side of the adjunction equality.
\item Fibre products of sets and of modules and elementwise verification of inverse constructions. used at proof lines 499-510. Establishes that the two constructions are mutually inverse.
\end{itemize}\medskip
\ifdefined\hypertarget\hypertarget{sourcepointer-1}{}\fi
\noindent\textit{Source reference.} Original lemma 0D2F: focused-seeds/AGP/sources/stacks-snapshot-20261011/tags.txt, source lines 16380--16380. The lemma's own canonical tag, quoted so that readers can resolve the label inside the delivered text; it points at the counted statement itself and carries no proof content.\par\medskip
\ifdefined\hypertarget\hypertarget{sourcepointer-2}{}\fi
\noindent\textit{Source reference.} Original lemma 0D2E: focused-seeds/AGP/sources/stacks-snapshot-20261011/tags.txt, source lines 16379--16379. Navigation pointer to a statement the counted proof refers to but does not contain; nothing load-bearing is delegated to it.\par\medskip
\ifdefined\hypertarget\hypertarget{sourcepointer-3}{}\fi
\noindent\textit{Source reference.} Original lemma 05DQ: focused-seeds/AGP/sources/stacks-snapshot-20261011/tags.txt, source lines 6963--6963. Navigation pointer to a statement the counted proof refers to but does not contain; nothing load-bearing is delegated to it.\par\medskip
\begin{lemma}
\label{lemma-modules}
Given a commutative diagram of rings
$$
\xymatrix{
R &
R' \ar[l] \\
B \ar[u] &
B' \ar[u] \ar[l]
}
$$
the functor (\ref{equation-modules}) has a right adjoint,
namely the functor
$$
F : (N, M', \varphi) \longmapsto N \times_\varphi M'
$$
(see proof for elucidation).
\end{lemma}

\begin{proof}
Given an object $(N, M', \varphi)$ of the category
$\text{Mod}_B \times_{\text{Mod}_R} \text{Mod}_{R'}$
we set
$$
N \times_\varphi M' = \{(n, m') \in N \times M' \mid
\varphi(n \otimes 1) = m' \otimes 1\text{ in }M' \otimes_{R'} R\}
$$
viewed as a $B'$-module.
The adjointness statement is that for a $B'$-module $L'$ and
a triple $(N, M', \varphi)$ we have
$$
\Hom_{B'}(L', N \times_\varphi M') =
\Hom_B(L' \otimes_{B'} B, N)
\times_{\Hom_R(L' \otimes_{B'} R, M' \otimes_{R'} R)}
\Hom_{R'}(L' \otimes_{B'} R', M')
$$
By Algebra, Lemma \ref{algebra-lemma-adjoint-tensor-restrict}
the right hand side is equal to
$$
\Hom_{B'}(L', N) \times_{\Hom_{B'}(L', M' \otimes_{R'} R)} \Hom_{B'}(L', M')
$$
Thus it is clear that for a pair $(g, f')$ of elements
of this fibre product we get an $B'$-linear map
$L' \to N \times_\varphi M'$, $l' \mapsto (g(l'), f'(l'))$.
Conversely, given a $B'$ linear map $g' : L' \to N \times_\varphi M'$
we can set $g$ equal to the composition $L' \to N \times_\varphi M' \to N$
and $f'$ equal to the composition $L' \to N \times_\varphi M' \to M'$.
These constructions are mutually inverse to each other and
define the desired isomorphism.
\end{proof}

\end{document}
